% TEMPLATE: quadrant / regime plot (where a rule holds in a plane).
% Two axes stated as full clauses, four cells, exactly one of them shaded and
% named; every cell carries a concrete instance, not a label.
\begin{figure}[H]
\centering
\pdfigurehue{pdcobalt}%
\begin{tikzpicture}[pd figure,x=1cm,y=1cm]
  \def\w{4.55}  \def\h{2.35}  \def\ox{2.35}  \def\oy{0.00}
  % the one regime the chapter is about, shaded and edged
  \draw[pd focus fill] ({\ox+\w},{\oy+\h}) rectangle ({\ox+2*\w},{\oy+2*\h});
  \draw[pd rule] (\ox,\oy) rectangle ({\ox+2*\w},{\oy+2*\h});
  \draw[pd rule] ({\ox+\w},\oy) -- ({\ox+\w},{\oy+2*\h});
  \draw[pd rule] (\ox,{\oy+\h}) -- ({\ox+2*\w},{\oy+\h});
  \foreach \cx/\cy/\head/\inst in {%
    0/1/{detect only}/{a plan stated in the model's context},
    1/1/{preventable}/{\texttt{git push}: net egress, spawn child},
    0/0/{detect only}/{a thought},
    1/0/{detect only}/{an effect whose guard reads unwitnessed state}}{
      \node[pd title,anchor=north] at ({\ox+\cx*\w+\w/2},{\oy+\cy*\h+\h-.28}) {\head};
      \node[pd label,text width=3.9cm,anchor=north] at ({\ox+\cx*\w+\w/2},{\oy+\cy*\h+\h-.78}) {\inst};}
  % axes stated as clauses, not as words
  \node[pd label,rotate=90,anchor=south,text width=4.3cm] at ({\ox-1.35},{\oy+\h})
    {the daemon controls the event};
  \node[pd label,anchor=north,text width=9cm] at ({\ox+\w},{\oy-.22})
    {the trigger is witnessed at the mediation boundary};
  \foreach \cx/\l in {0/no, 1/yes}{
    \node[pd label,anchor=south] at ({\ox+\cx*\w+\w/2},{\oy+2*\h+.08}) {\l};}
  \foreach \cy/\l in {0/no, 1/yes}{
    \node[pd label,anchor=east] at ({\ox-.20},{\oy+\cy*\h+\h/2}) {\l};}
\end{tikzpicture}
\caption{Preventability is one cell, not a spectrum. An event is preventable
exactly when it is controllable by the daemon (top row) \emph{and} its trigger
is witnessed at the mediation boundary (right column): the shaded cell. Each of
the other three cells names a concrete case that is detect-only, so the reader
can test the rule against an instance rather than a label. [internal]}
\label{fig:tpl-quadrant}
\end{figure}
